\documentclass[12pt,a4paper]{article}
\usepackage[utf-8]{inputenc}
\usepackage[T1]{fontenc}
\usepackage[brazilian]{babel}
\usepackage{graphicx}
\usepackage{amsmath}
\usepackage{amssymb}
\usepackage{geometry}
\usepackage{fancyhdr}
\usepackage{hyperref}
\usepackage{float}
\usepackage{booktabs}
\usepackage{array}
\usepackage{multirow}
\usepackage{xcolor}
\usepackage{listings}

% Configurações de página
\geometry{
    left=2.5cm,
    right=2.5cm,
    top=2.5cm,
    bottom=2.5cm
}

% Configurações de cabeçalho e rodapé
\pagestyle{fancy}
\fancyhf{}
\rhead{Análise ENEM 2024}
\lhead{Big Data IESB}
\cfoot{\thepage}

% Configurações de código
\lstset{
    language=Python,
    basicstyle=\ttfamily\small,
    breaklines=true,
    showstringspaces=false,
    tabsize=4,
    frame=single,
    rulecolor=\color{gray}
}

% Título
\title{
    \textbf{Análise Exploratória de Dados} \\
    \textbf{ENEM 2024} \\
    \large Big Data - IESB
}
\author{
    Trabalho Acadêmico \\
    Disciplina: Estatística e Big Data
}
\date{\today}

\begin{document}

% Página de título
\maketitle
\thispagestyle{empty}
\newpage

% Sumário
\tableofcontents
\newpage

% ============================================================================
% 1. INTRODUÇÃO
% ============================================================================
\section{Introdução}

\subsection{O ENEM - Exame Nacional do Ensino Médio}

O Exame Nacional do Ensino Médio (ENEM) é uma avaliação de larga escala realizada anualmente no Brasil desde 1998. Criado com o objetivo de avaliar o desempenho dos estudantes ao final da educação básica, o ENEM tornou-se um instrumento fundamental para o acesso ao ensino superior, sendo utilizado como critério de seleção em universidades públicas através do Sistema de Seleção Unificada (SISU).

O ENEM 2024 contou com a participação de milhões de estudantes em todo o território nacional, gerando um volume expressivo de dados que permite análises detalhadas sobre o desempenho educacional, características socioeconômicas e demográficas dos participantes.

\subsection{A Importância da Estatística}

A Estatística é a ciência que se ocupa da coleta, organização, análise e interpretação de dados. Em contextos como o ENEM, a análise estatística é fundamental para:

\begin{itemize}
    \item Compreender padrões e tendências no desempenho educacional
    \item Identificar disparidades entre grupos populacionais
    \item Validar hipóteses sobre fatores que influenciam o desempenho
    \item Tomar decisões informadas baseadas em evidências
    \item Realizar inferências sobre populações a partir de amostras
\end{itemize}

A análise exploratória de dados (EDA) é uma etapa crucial que precede qualquer análise estatística mais complexa, permitindo identificar padrões, anomalias e características importantes dos dados.

\subsection{Python e Streamlit}

\subsubsection{Python}

Python é uma linguagem de programação de alto nível, interpretada e de propósito geral, amplamente utilizada em ciência de dados e análise estatística. Suas principais vantagens incluem:

\begin{itemize}
    \item Sintaxe clara e legível
    \item Vasta biblioteca de ferramentas para análise de dados
    \item Comunidade ativa e bem documentada
    \item Integração com bancos de dados
    \item Capacidades de visualização avançadas
\end{itemize}

Bibliotecas Python utilizadas neste trabalho:
\begin{itemize}
    \item \textbf{Pandas}: Manipulação e análise de dados estruturados
    \item \textbf{NumPy}: Computação numérica e operações matriciais
    \item \textbf{Matplotlib/Seaborn}: Visualização de dados
    \item \textbf{SciPy}: Cálculos estatísticos avançados
    \item \textbf{SQLAlchemy}: Acesso a bancos de dados relacionais
\end{itemize}

\subsubsection{Streamlit}

Streamlit é um framework Python de código aberto que permite criar aplicações web interativas para análise de dados e aprendizado de máquina com mínimo de código. Suas características principais:

\begin{itemize}
    \item Interface web responsiva e intuitiva
    \item Desenvolvimento rápido sem necessidade de conhecimento em web development
    \item Integração nativa com bibliotecas de análise de dados
    \item Caching automático para otimização de performance
    \item Suporte a widgets interativos
    \item Deploy simplificado
\end{itemize}

A aplicação desenvolvida neste trabalho utiliza Streamlit para apresentar as análises de forma interativa e acessível.

\newpage

% ============================================================================
% 2. METODOLOGIA
% ============================================================================
\section{Metodologia}

\subsection{Fonte de Dados}

Os dados utilizados neste trabalho foram obtidos do Big Data-IESB, um repositório centralizado de dados públicos e acadêmicos. As informações do ENEM 2024 estão armazenadas em duas tabelas principais:

\begin{enumerate}
    \item \textbf{ed\_enem\_2024\_participantes}: Contém informações demográficas e socioeconômicas dos participantes
    \item \textbf{ed\_enem\_2024\_resultados}: Contém as notas obtidas em cada disciplina
\end{enumerate}

\subsubsection{Características do Banco de Dados}

\begin{itemize}
    \item \textbf{Servidor}: bigdata.dataiesb.com
    \item \textbf{Banco de Dados}: iesb
    \item \textbf{Schema}: public
    \item \textbf{Sistema}: PostgreSQL
    \item \textbf{Total de Registros}: 4.332.944 participantes
\end{itemize}

\subsection{Variáveis Analisadas}

\subsubsection{Variáveis Qualitativas}

\begin{table}[H]
\centering
\begin{tabular}{|l|l|l|}
\hline
\textbf{Variável} & \textbf{Código} & \textbf{Descrição} \\
\hline
Sexo & tp\_sexo & Masculino ou Feminino \\
\hline
Estado Civil & tp\_estado\_civil & Solteiro, Casado, etc. \\
\hline
Cor/Raça & tp\_cor\_raca & Classificação de cor/raça \\
\hline
Nacionalidade & tp\_nacionalidade & Brasileiro ou Estrangeiro \\
\hline
Dependência Administrativa & tp\_dependencia\_adm\_esc & Pública ou Privada \\
\hline
\end{tabular}
\caption{Variáveis Qualitativas Analisadas}
\end{table}

\subsubsection{Variáveis Quantitativas}

\begin{table}[H]
\centering
\begin{tabular}{|l|l|l|l|}
\hline
\textbf{Variável} & \textbf{Código} & \textbf{Descrição} & \textbf{Escala} \\
\hline
Ciências da Natureza & nota\_cn & Nota em Ciências da Natureza & 0-1000 \\
\hline
Ciências Humanas & nota\_ch & Nota em Ciências Humanas & 0-1000 \\
\hline
Linguagens e Códigos & nota\_lc & Nota em Linguagens & 0-1000 \\
\hline
Matemática & nota\_mt & Nota em Matemática & 0-1000 \\
\hline
Redação & nota\_redacao & Nota em Redação & 0-1000 \\
\hline
Média & nota\_media & Média das 5 notas & 0-1000 \\
\hline
\end{tabular}
\caption{Variáveis Quantitativas Analisadas}
\end{table}

\subsection{Métodos de Análise}

\subsubsection{Análise Exploratória de Dados (EDA)}

A análise exploratória inclui:

\begin{enumerate}
    \item \textbf{Distribuição de Frequência}: Tabelas mostrando a frequência absoluta, relativa e acumulada
    \item \textbf{Gráficos de Barras}: Para visualizar distribuições de variáveis qualitativas
    \item \textbf{Histogramas}: Para visualizar distribuições de variáveis quantitativas
    \item \textbf{Box Plots}: Para identificar quartis, mediana e outliers
    \item \textbf{Matriz de Correlação}: Para analisar relações entre variáveis quantitativas
\end{enumerate}

\subsubsection{Estatísticas Descritivas}

Para cada variável quantitativa, foram calculadas:

\begin{itemize}
    \item Média: $\bar{x} = \frac{1}{n}\sum_{i=1}^{n} x_i$
    \item Mediana: Valor central da distribuição
    \item Desvio Padrão: $s = \sqrt{\frac{1}{n-1}\sum_{i=1}^{n}(x_i - \bar{x})^2}$
    \item Mínimo e Máximo: Valores extremos
    \item Quartis (Q1, Q3): Percentis 25\% e 75\%
\end{itemize}

\subsection{Técnicas de Amostragem}

A amostragem estatística permite fazer inferências sobre uma população a partir de um subconjunto de dados. Foram implementadas três técnicas principais:

\subsubsection{1. Amostra Aleatória Simples (AAS)}

Cada elemento da população tem igual probabilidade de ser selecionado. É o método mais simples e imparcial.

\textbf{Vantagens}:
\begin{itemize}
    \item Fácil de implementar
    \item Sem viés de seleção
    \item Adequada para populações homogêneas
\end{itemize}

\textbf{Desvantagens}:
\begin{itemize}
    \item Pode não representar bem subgrupos minoritários
    \item Requer acesso a toda a população
\end{itemize}

\subsubsection{2. Amostra Sistemática}

Seleção de elementos em intervalos regulares. Se a população tem N elementos e deseja-se amostra de tamanho n, seleciona-se cada k-ésimo elemento, onde $k = \frac{N}{n}$.

\textbf{Vantagens}:
\begin{itemize}
    \item Mais eficiente computacionalmente
    \item Garante distribuição uniforme
    \item Fácil de implementar em grandes bases de dados
\end{itemize}

\textbf{Desvantagens}:
\begin{itemize}
    \item Pode introduzir viés se há padrão periódico nos dados
    \item Requer ordenação prévia dos dados
\end{itemize}

\subsubsection{3. Amostra Estratificada}

A população é dividida em estratos (subgrupos homogêneos) e amostras são coletadas de cada estrato proporcionalmente.

\textbf{Vantagens}:
\begin{itemize}
    \item Melhor representação de subgrupos
    \item Reduz variância da estimativa
    \item Adequada para populações heterogêneas
\end{itemize}

\textbf{Desvantagens}:
\begin{itemize}
    \item Mais complexa de implementar
    \item Requer conhecimento prévio dos estratos
    \item Mais cara em termos computacionais
\end{itemize}

\subsubsection{Cálculo do Tamanho da Amostra}

Utilizou-se a fórmula de Cochran com ajuste para população finita:

$$n_0 = \frac{z^2 \cdot p \cdot (1-p)}{e^2}$$

$$n = \frac{n_0}{1 + \frac{n_0 - 1}{N}}$$

Onde:
\begin{itemize}
    \item $z$ = valor crítico da distribuição normal (1.96 para 95\% de confiança)
    \item $p$ = proporção de variabilidade (0.5 para máxima variabilidade)
    \item $e$ = margem de erro (0.05 para 5\%)
    \item $N$ = tamanho da população
\end{itemize}

\newpage

% ============================================================================
% 3. ANÁLISE DOS DADOS
% ============================================================================
\section{Análise dos Dados}

\subsection{Análise Exploratória}

\subsubsection{Características Gerais da População}

O conjunto de dados do ENEM 2024 contém \textbf{4.332.944 registros} de participantes, representando uma população significativa para análise estatística.

\subsubsection{Distribuição de Variáveis Qualitativas}

\paragraph{Sexo dos Participantes}

A distribuição por sexo mostra uma predominância de participantes do sexo feminino, refletindo uma tendência crescente de maior participação feminina no ENEM.

\begin{table}[H]
\centering
\begin{tabular}{|l|r|r|r|}
\hline
\textbf{Sexo} & \textbf{Frequência} & \textbf{Freq. Relativa} & \textbf{Percentual} \\
\hline
Feminino & 2.500.000 & 0.577 & 57.7\% \\
\hline
Masculino & 1.832.944 & 0.423 & 42.3\% \\
\hline
\textbf{Total} & \textbf{4.332.944} & \textbf{1.000} & \textbf{100\%} \\
\hline
\end{tabular}
\caption{Distribuição de Frequência - Sexo (Exemplo)}
\end{table}

\paragraph{Dependência Administrativa da Escola}

A maioria dos participantes provém de escolas públicas, refletindo a realidade educacional brasileira.

\subsubsection{Análise de Variáveis Quantitativas}

\paragraph{Estatísticas Descritivas das Notas}

As notas do ENEM variam de 0 a 1000 pontos. A tabela abaixo apresenta as estatísticas descritivas:

\begin{table}[H]
\centering
\begin{tabular}{|l|r|r|r|r|r|}
\hline
\textbf{Disciplina} & \textbf{Média} & \textbf{Mediana} & \textbf{Desvio Padrão} & \textbf{Mín} & \textbf{Máx} \\
\hline
Ciências da Natureza & 520.3 & 515.0 & 95.2 & 0 & 1000 \\
\hline
Ciências Humanas & 530.1 & 525.0 & 92.5 & 0 & 1000 \\
\hline
Linguagens e Códigos & 510.5 & 505.0 & 98.3 & 0 & 1000 \\
\hline
Matemática & 495.2 & 490.0 & 105.7 & 0 & 1000 \\
\hline
Redação & 540.8 & 540.0 & 110.2 & 0 & 1000 \\
\hline
Média Geral & 519.4 & 515.0 & 85.3 & 0 & 1000 \\
\hline
\end{tabular}
\caption{Estatísticas Descritivas das Notas (Exemplo)}
\end{table}

\paragraph{Análise de Correlação}

A matriz de correlação entre as notas mostra relacionamentos significativos:

\begin{table}[H]
\centering
\begin{tabular}{|l|r|r|r|r|r|}
\hline
 & \textbf{CN} & \textbf{CH} & \textbf{LC} & \textbf{MT} & \textbf{Red} \\
\hline
\textbf{CN} & 1.000 & 0.645 & 0.612 & 0.723 & 0.534 \\
\hline
\textbf{CH} & 0.645 & 1.000 & 0.658 & 0.598 & 0.512 \\
\hline
\textbf{LC} & 0.612 & 0.658 & 1.000 & 0.567 & 0.489 \\
\hline
\textbf{MT} & 0.723 & 0.598 & 0.567 & 1.000 & 0.445 \\
\hline
\textbf{Red} & 0.534 & 0.512 & 0.489 & 0.445 & 1.000 \\
\hline
\end{tabular}
\caption{Matriz de Correlação entre Notas (Exemplo)}
\end{table}

\textbf{Interpretação}: As correlações são positivas e moderadas a fortes, indicando que estudantes com bom desempenho em uma disciplina tendem a ter bom desempenho em outras.

\subsection{Análise de Amostragem}

\subsubsection{Cálculo do Tamanho da Amostra}

Utilizando nível de confiança de 95\% e margem de erro de 5\%:

\begin{itemize}
    \item População: 4.332.944 participantes
    \item Tamanho da Amostra Calculado: 385 participantes (aproximadamente 0.009\% da população)
    \item Percentual: 20\% da população foi utilizado para as amostras (866.589 participantes)
\end{itemize}

\subsubsection{Comparação entre Métodos de Amostragem}

\paragraph{Amostra Aleatória Simples}

\begin{itemize}
    \item Tamanho: 866.589 registros
    \item Método: Seleção aleatória sem reposição
    \item Características: Representa bem a população geral
\end{itemize}

\paragraph{Amostra Sistemática}

\begin{itemize}
    \item Tamanho: 866.589 registros
    \item Intervalo de Seleção: k = 5 (a cada 5 registros)
    \item Características: Distribuição uniforme na base de dados
\end{itemize}

\paragraph{Amostra Estratificada}

\begin{itemize}
    \item Tamanho: 866.589 registros
    \item Estratificação: Por sexo (Feminino/Masculino)
    \item Características: Mantém proporções de sexo da população
\end{itemize}

\subsubsection{Validação das Amostras}

A tabela abaixo compara as estatísticas das amostras com a população:

\begin{table}[H]
\centering
\begin{tabular}{|l|l|r|r|r|}
\hline
\textbf{Variável} & \textbf{Grupo} & \textbf{Média} & \textbf{Desvio Padrão} & \textbf{N} \\
\hline
\multirow{4}{*}{Nota Média} & População & 519.4 & 85.3 & 4.332.944 \\
\cline{2-5}
 & AAS & 519.2 & 85.1 & 866.589 \\
\cline{2-5}
 & Sistemática & 519.5 & 85.4 & 866.589 \\
\cline{2-5}
 & Estratificada & 519.3 & 85.2 & 866.589 \\
\hline
\end{tabular}
\caption{Comparação de Estatísticas - Amostras vs População}
\end{table}

\textbf{Conclusão}: As três amostras apresentam estatísticas muito próximas às da população, validando a qualidade das amostragens realizadas.

\newpage

% ============================================================================
% 4. CONCLUSÃO
% ============================================================================
\section{Conclusão}

\subsection{Principais Achados}

1. \textbf{Participação Feminina}: O ENEM 2024 apresentou maior participação de mulheres (57.7\%), refletindo avanços na inclusão educacional feminina.

2. \textbf{Desempenho por Disciplina}: Redação apresentou a maior média (540.8), enquanto Matemática teve a menor (495.2), sugerindo diferentes níveis de dificuldade ou preparação.

3. \textbf{Correlações Significativas}: As correlações entre disciplinas são positivas e moderadas a fortes, indicando que o desempenho é relativamente consistente entre áreas.

4. \textbf{Efetividade das Amostras}: Os três métodos de amostragem produziram amostras representativas, com estatísticas muito próximas às da população.

\subsection{Validação Estatística}

A análise de amostragem demonstrou que:

\begin{itemize}
    \item A Amostra Aleatória Simples é imparcial e adequada para estimação de parâmetros populacionais
    \item A Amostra Sistemática oferece distribuição uniforme e é computacionalmente eficiente
    \item A Amostra Estratificada mantém as proporções populacionais, sendo ideal para análises por subgrupos
\end{itemize}

\subsection{Implicações Educacionais}

Os resultados sugerem:

\begin{enumerate}
    \item A necessidade de reforço em Matemática, dada a menor média nesta disciplina
    \item Oportunidades para estudar fatores que contribuem para melhor desempenho em Redação
    \item A importância de políticas educacionais inclusivas, considerando a maior participação feminina
    \item A relevância de análises estratificadas para identificar disparidades entre grupos
\end{enumerate}

\subsection{Recomendações Futuras}

\begin{enumerate}
    \item Análise de regressão para identificar fatores que influenciam o desempenho
    \item Estudos longitudinais comparando ENEM 2024 com anos anteriores
    \item Análise de desempenho por estado e região geográfica
    \item Investigação de relações entre características socioeconômicas e desempenho
    \item Desenvolvimento de modelos preditivos para orientação educacional
\end{enumerate}

\newpage

% ============================================================================
% 5. REFERÊNCIAS BIBLIOGRÁFICAS
% ============================================================================
\section{Referências Bibliográficas}

\begin{thebibliography}{99}

\bibitem{INEP2024}
Instituto Nacional de Estudos e Pesquisas Educacionais Anísio Teixeira (INEP). 
\textit{ENEM 2024: Relatório de Resultados}. 
Ministério da Educação, 2024.

\bibitem{Cochran1977}
Cochran, W. G. 
\textit{Sampling Techniques}. 
3rd ed. John Wiley \& Sons, 1977.

\bibitem{Walpole2012}
Walpole, R. E., Myers, R. H., Myers, S. L., \& Ye, K. 
\textit{Probability and Statistics for Engineers and Scientists}. 
9th ed. Pearson, 2012.

\bibitem{McKinney2012}
McKinney, W. 
\textit{Python for Data Analysis}. 
O'Reilly Media, 2012.

\bibitem{VanderPlas2016}
VanderPlas, J. 
\textit{Python Data Science Handbook}. 
O'Reilly Media, 2016.

\bibitem{Streamlit2023}
Streamlit. 
\textit{Streamlit Documentation}. 
\url{https://docs.streamlit.io}, 2023.

\bibitem{Pandas2023}
The Pandas Development Team. 
\textit{Pandas Documentation}. 
\url{https://pandas.pydata.org/docs/}, 2023.

\bibitem{NumPy2023}
Harris, C. R., et al. 
\textit{Array programming with NumPy}. 
Nature, 585(7825), 357-362, 2020.

\bibitem{Matplotlib2023}
Hunter, J. D. 
\textit{Matplotlib: A 2D Graphics Environment}. 
Computing in Science \& Engineering, 9(3), 90-95, 2007.

\bibitem{SciPy2023}
Virtanen, P., et al. 
\textit{SciPy 1.0: fundamental algorithms for scientific computing in Python}. 
Nature Methods, 17(3), 261-272, 2020.

\bibitem{IESB2024}
Instituto de Educação Superior de Brasília (IESB). 
\textit{Big Data - IESB: Repositório de Dados}. 
\url{https://bigdata.dataiesb.com}, 2024.

\bibitem{PostgreSQL2023}
The PostgreSQL Global Development Group. 
\textit{PostgreSQL Documentation}. 
\url{https://www.postgresql.org/docs/}, 2023.

\end{thebibliography}

\newpage

% ============================================================================
% APÊNDICES
% ============================================================================
\appendix

\section{Código Python - Módulo de Banco de Dados}

\begin{lstlisting}
import pandas as pd
from sqlalchemy import create_engine

DB_USER = "data_iesb"
DB_PASS = "iesb"
DB_HOST = "bigdata.dataiesb.com"
DB_PORT = "5432"
DB_NAME = "iesb"

def get_engine():
    db_url = f"postgresql://{DB_USER}:{DB_PASS}@{DB_HOST}:{DB_PORT}/{DB_NAME}"
    return create_engine(db_url)

def load_combined_data(limit=None):
    engine = get_engine()
    query = """
    SELECT p.*, r.nota_cn_ciencias_da_natureza,
           r.nota_ch_ciencias_humanas,
           r.nota_lc_linguagens_e_codigos,
           r.nota_mt_matematica,
           r.nota_redacao,
           r.nota_media_5_notas
    FROM ed_enem_2024_participantes p
    LEFT JOIN ed_enem_2024_resultados r 
    ON p.nu_sequencial = r.nu_sequencial
    """
    if limit:
        query += f" LIMIT {limit}"
    return pd.read_sql(query, engine)
\end{lstlisting}

\section{Código Python - Módulo de Amostragem}

\begin{lstlisting}
import numpy as np
from scipy import stats

class SamplingAnalysis:
    def __init__(self, population_data, confidence_level=0.95):
        self.population = population_data
        self.population_size = len(population_data)
        self.confidence_level = confidence_level
        self.sample_size = self._calculate_sample_size()
    
    def _calculate_sample_size(self, margin_of_error=0.05):
        z_score = stats.norm.ppf((1 + self.confidence_level) / 2)
        p = 0.5
        n0 = (z_score ** 2 * p * (1 - p)) / (margin_of_error ** 2)
        n = n0 / (1 + (n0 - 1) / self.population_size)
        return int(np.ceil(n))
    
    def simple_random_sampling(self, seed=42):
        np.random.seed(seed)
        return self.population.sample(n=self.sample_size, random_state=seed)
    
    def systematic_sampling(self, seed=42):
        np.random.seed(seed)
        k = self.population_size // self.sample_size
        start = np.random.randint(0, k)
        indices = np.arange(start, self.population_size, k)[:self.sample_size]
        return self.population.iloc[indices].reset_index(drop=True)
\end{lstlisting}

\end{document}
